% Generated draft: requires mathematical and source review.
\documentclass[11pt]{article}
\usepackage[margin=25mm]{geometry}
\usepackage{fontspec}
\setmainfont{Latin Modern Roman}
\usepackage{amsmath,amssymb,mathtools,unicode-math}
\setmathfont{Latin Modern Math}
\usepackage{enumitem,needspace,xurl,hyperref}
\setlength{\parindent}{0pt}
\setlength{\parskip}{5pt}
\setlength{\emergencystretch}{3em}
\newcommand{\N}{\mathbb{N}}
\newcommand{\Z}{\mathbb{Z}}
\newcommand{\Q}{\mathbb{Q}}
\newcommand{\R}{\mathbb{R}}
\newcommand{\C}{\mathbb{C}}
\newcommand{\F}{\mathbb{F}}
\newcommand{\A}{\mathbb{A}}
\newcommand{\PP}{\mathbb{P}}
\newcommand{\E}{\mathbb{E}}
\newcommand{\eps}{\varepsilon}
\newcommand{\dd}{\,\mathrm{d}}
\begin{document}
\section*{UnsolvedMath ID: 30006248}
\textbf{Source problem code:} catalogue entry
\subsection{Definitions and mathematical statement}
For a finite group \(G\), let \(Z(G)\) be its centre and let \(\operatorname{Irr}(G)\) be the set of irreducible complex characters of \(G\). For \(\chi\in\operatorname{Irr}(G)\), its degree is \(\chi(1)\).

The apparent intended definition of \emph{almost quasisimple} in the supplied record is that, for some finite nonabelian simple group \(S\),
\[
S\trianglelefteq G/Z(G)\leq\operatorname{Aut}(S),
\]
where \(S\) is identified with its inner automorphism group and the inclusions are understood through an identification of \(G/Z(G)\) with a subgroup of \(\operatorname{Aut}(S)\). This reading requires an explicit correction: the original text says ``some finite non-abelian simple group \(G\)'' although its displayed condition uses \(S\). The intended definition requires source confirmation.

\textbf{Requested statement, conditional on that interpretation.} For every finite almost quasisimple group \(G\) and every \(g\in G\setminus Z(G)\), determine an explicit real exponent \(\alpha_G(g)\) with
\[
0\leq\alpha_G(g)<1
\quad\text{and}\quad
|\chi(g)|\leq\chi(1)^{\alpha_G(g)}
\qquad\text{for every }\chi\in\operatorname{Irr}(G),
\]
and make this exponent as small as possible. The subscript \(G\) makes the ambient group explicit; the source writes \(\alpha(g)\). The same exponent must work for all irreducible characters of the fixed group and element. No restriction to faithful characters is supplied.

The existence assertion needs mathematical review in the stated generality. In particular, for every character of degree greater than one, the requested inequality forces \(|\chi(g)|<\chi(1)\). Thus a noncentral element in the kernel of any nonlinear irreducible representation would obstruct the requested exponent. The source does not explain how such obstructions, including those potentially arising from characters inflated from quotients, are excluded. No missing hypothesis is imposed here. The source also does not specify the required form of an ``explicit'' exponent or whether ``as small as possible'' requires an attained minimum or an optimal estimate within a specified family.

\textbf{Supplied contextual results, not independently verified.} The excerpt attributes to its Theorem 2, reference [6], a uniform bound for every finite quasisimple group \(G\) of Lie type, every \(\chi\in\operatorname{Irr}(G)\), and every \(g\in G\). The apparent reading of its damaged display is that there is an explicit absolute constant \(c>0\) such that
\[
|\chi(g)|\leq\chi(1)^{\,1-c\frac{\log|g^G|}{\log|G|}},
\qquad
g^G=\{x^{-1}gx:x\in G\}.
\]
This is a typographical reconstruction requiring verification. Its supplied scope is quasisimple groups of Lie type, rather than the entire class in the requested statement. The excerpt says it builds on references [2, 4, 5, 8], whose bibliographic entries are absent.

The excerpt next considers classical groups \(G=\operatorname{Cl}(V)=\operatorname{Cl}_n(q)\), with \(V=\mathbb F_q^n\). It describes the support of \(g\) as the least codimension of an eigenspace after scalar extension and says that \(\operatorname{supp}(g)/n\) and \(\log|g^G|/\log|G|\) have the same magnitude. The scalar field and scalar extension in the extracted support definition are damaged and cannot be fixed unambiguously from this record. The notion of character level is referred to [4, 5] but is not defined.

Theorem 3, attributed to [7], is reported with \(n,m,j\in\mathbb Z_{\geq1}\), a prime power \(q\), a sign \(\varepsilon\in\{+,-\}\), and
\[
n\geq8j^2+4j+4mj+3.
\]
It concerns an irreducible complex character \(\chi\) of level \(j\) of \(G=\operatorname{GL}_n(q)\), \(\operatorname{GU}_n(q)\), or \(\operatorname{SO}^{\varepsilon}_n(q)\). An additional condition in the orthogonal case appears to be \(2\nmid q\), but the extraction is damaged. For elements of support \(m\), the apparent estimate is that there exists a root of unity \(\beta\in\mathbb C\) such that
\[
\left|q^{mj}\frac{\chi(g)}{\chi(1)}-\beta\right|<q^{-n/4}.
\]
The display, the support and level conventions, and the dependence of \(\beta\) require verification; this reconstruction is not asserted as a verified theorem.

The excerpt further reports an application of Theorems 2 and 3 to an asymptotic version of Thompson's conjecture. It states the full conjecture as follows: every finite nonabelian simple group \(H\) has a conjugacy class \(C\subseteq H\) with \(C^2=H\), where \(C^2=\{xy:x,y\in C\}\). The precise asymptotic range is not supplied. This contextual conjecture is distinct from the requested character-bound problem.
\subsection{Short English statement}
For each finite almost quasisimple group and each noncentral element, seek an explicit exponent below one that bounds every irreducible character value by that power of its degree, with the exponent as small as possible. The supplied definition has a variable error, and the unrestricted quantifier over characters may require additional hypotheses. This is a statement-only draft requiring human mathematical review.
\subsection{Sources}
\begin{itemize}
\item Local UnsolvedMath registry, permanent ID \(30006248\), source problem code \texttt{catalogue entry}, title ``Optimal Character Bounds for Almost Quasisimple Groups,'' in the supplied Oberwolfach Reports Open Problems collection. The supplied record labels the problem \texttt{partially\_solved} and its statement \texttt{corrected\_verified}; these are registry labels, not independent verification. Its verification note says that only formula delimiters, notation, grammar, and typography were repaired.
\item \emph{Finite Groups, Fusion Systems and Applications} (2025), Oberwolfach Reports, pp. 689--752, DOI 10.4171/owr/2025/16: \url{https://doi.org/10.4171/owr/2025/16}. This is the supplied report citation and the literature-source link labelled ``problem and current results.'' The excerpt identifies the contribution as ``Character Estimates for Finite Classical Groups and Thompson's Conjecture,'' by Pham Huu Tiep, joint work with Michael Larsen. The registry supplies no separate proposer attribution.
\item Supplied literature source, ``Larsen--Tiep, uniform character bounds'': \url{https://annals.math.princeton.edu/2024/200-1}. The supplied link identifies a journal volume/issue rather than a complete individual article citation.
\item Supplied literature source, ``Larsen--Tiep, classical estimates'': \url{https://arxiv.org/abs/2403.09047}.
\item Literature assessment checked on \textbf{2026-08-22}: this is supplied dated curation, not independently verified. It reports proved uniform exponential character bounds and sharp estimates in important classical-group regimes, while describing optimal explicit exponents for every almost quasisimple group and element as still open. The background explicitly calls this dated literature triage derived from the supplied OWR curation and asks for independent verification of time-sensitive status and cited claims.
\item The excerpt cites numbered references [2, 4, 5, 6, 7, 8] without their bibliography. Their exact identities and their correspondence with the supplied links cannot be established from this record alone. It attributes the claim that Thompson's conjecture first appeared in Kourovka's Notebook to a private communication by Khukhro; no further citation is supplied.
\end{itemize}
\end{document}
